\def\stat{bazilevskii}





\def\tit{ПРОГРАММА ПОСТРОЕНИЯ ВПОЛНЕ ИНТЕРПРЕТИРУЕМЫХ И~RTF-АДЕКВАТНЫХ 


ЛИНЕЙНЫХ РЕГРЕССИОННЫХ МОДЕЛЕЙ}





\def\titkol{Программа построения вполне интерпретируемых и~RTF-адекватных 


линейных %регрессионных 


моделей}





\def\aut{М.\,П.~Базилевский$^1$}





\def\autkol{М.\,П.~Базилевский}





\titel{\tit}{\aut}{\autkol}{\titkol}





\index{Базилевский М.\,П.}


\index{Bazilevskiy M.\,P.}





%{\renewcommand{\thefootnote}{\fnsymbol{footnote}}


%\footnotetext[1]


%{Работа частично поддержана РФФИ (проект 18-29-03124-мк).}} 





\renewcommand{\thefootnote}{\arabic{footnote}}


\footnotetext[1]{Иркутский государственный университет путей сообщения, кафедра математики, 


\mbox{mik2178@yandex.ru}}





 \label{st\stat}





 \thispagestyle{headings}


 


 \vspace*{-12pt} 











  \Abst{Статья посвящена проблеме отбора <<информативных>> регрессоров (ОИР)


  в~регрессионных моделях, оцениваемых с~помощью метода наименьших квадратов (МНК). 


Построенные в~результате такого отбора модели часто оказываются неадекватными и~плохо 


интерпретируемыми. В~работе впервые сформулированы определения <<вполне 


интерпретируемой>> и~<<RTF-адекват\-ной>> регрессионной модели. Рассмотрен ранее 


предложенный эффективный алгоритм решения задачи ОИР. 


На его основе разработан алгоритм построения вполне интерпретируемых 


и~RTF-адекват\-ных линейных регрессионных моделей. В~нем для каждой регрессии 


последовательно осуществляется проверка: <<информативности>> переменных, 


мультиколлинеарности, соответствия знаков коэффициентов физическому смыслу факторов, 


адекватности модели по коэффициенту детерминации и~значимости в~целом по  


F-кри\-те\-рию Фишера, а также значимости коэффициентов по t-кри\-те\-рию Стьюдента. 


Предложенный алгоритм реализован в~виде программы для эконометрического пакета Gretl. 


Разработанная программа универсальна и~может быть использована для решения широкого 


круга задач анализа данных.}


  


  \KW{отбор <<информативных>> регрессоров; метод наименьших квадратов; вполне 


интерпретируемая и~RTF-адекватная регрессия; критерий <<информативности>> 


переменных; мультиколлинеарность; F-кри\-те\-рий Фишера; t-кри\-те\-рий Стьюдента}





\DOI{10.14357\!/\!08696527210402} 





\vskip 10pt plus 9pt minus 6pt





 \thispagestyle{headings}


 


 %\vspace*{-20pt}





\section{Введение}





  Один из этапов регрессионного анализа~[1]~--- спецификация модели, т.\,е.\ 


выбор состава переменных и~математической формы связи между ними. 


Определение спецификации регрессии на практике зачастую сводится 


к~решению задачи ОИР, более 


известной в~зарубежной литературе как <<subset selection>>, <<variable 


selection>> или <<feature selection>> in regression~[2, 3]. В~отечественной 


литературе описание методов ОИР можно найти в~работе~[4]. Решение задачи 


ОИР может осуществляться не только по ка\-ко\-му-то одному критерию 


адекватности, а одновременно по нескольким критериям сразу. Такая 


технология многокритериального выбора получила название <<конкурс>> 


моделей~[5, 6].


  


  В работе~[7] автором был проведен анализ и~выявлены сле\-ду\-ющие 


недостатки программного комплекса автоматизации процесса построения 


регрессионных моделей (ПК АППРМ), предназначенного для проведения 


<<конкурса>> моделей.


  \begin{enumerate}[1.]


\item Полученные в~результате <<конкурса>> модели часто представляют 


собой сложные нелинейные уравнения, коэффициенты которых весьма 


затруднительно интерпретировать. При этом игнорируется проб\-ле\-ма 


мультиколлинеарности, которая может лишать смысла интерпретацию 


коэффициентов модели.


\item Полученные в~результате <<конкурса>> модели часто оказываются 


неадекватными по коэффициенту детерминации~$R^2$, незначимыми  


по F-кри\-те\-рию Фишера, а их оценки~--- незначимыми по t-кри\-те\-рию 


Стьюдента.


\end{enumerate}


  


  На основе выявленных недостатков в~работе~[7] было введено понятие 


<<хорошо интерпретируемая качественная регрессия>>, удовлетворяющая 


пяти условиям, и~разработан фундаментальный блок алгоритмов построения 


таких моделей.


  


  Цель данной работы~--- создание конкретного алгоритма построения хорошо 


интерпретируемых и~качественных линейных регрессий, а также его 


реализация в~виде программы.


  


\section{Вполне интерпретируемые и~RTF-адекватные регрессионные 


модели}





  Размышления над предложенным в~[7] понятием <<хорошо 


интерпретируемая качественная модель>> привели к~тому, что целесообразнее 


заменить этот термин следующими двумя определениями.


  \smallskip


  


  \noindent


  \textbf{Определение~1.}\ Регрессионная модель считается вполне 


интерпретируемой, если она удовлетворяет трем условиям:


  \begin{enumerate}[(1)]


\item ее спецификация изначально выбрана так, что после оценивания модели 


можно объяснить любой ее коэффициент или некоторый его аналог, за 


исключением, быть может, свободного члена;


\item знаки коэффициентов оцененной модели соответствуют физическому 


смыслу входящих в~уравнение факторов;


\item эффект мультиколлинеарности незначителен.


\end{enumerate}





  \noindent


  \textbf{Определение~2.}\  Регрессионная модель, оцененная с~помощью 


МНК, считается RTF-адекват\-ной, если она удовлетворяет трем условиям:


  \begin{enumerate}[(1)]


\item модель адекватна по коэффициенту детерминации~$R^2$;


\item коэффициенты модели значимы по t-кри\-те\-ри\-ям Стьюдента;


\item модель значима в~целом по F-кри\-те\-рию Фишера.


\end{enumerate}





\section{Эффективный алгоритм решения задачи отбора 


<<информативных>> регрессоров}





  Рассмотрим модель множественной линейной регрессии


  \begin{equation}


  y_i= \alpha_0+\alpha_1x_{i1} +\alpha_2x_{i2}+\cdots + \alpha_l 


x_{il}+\varepsilon_i\,,\enskip i=\overline{1,n}\,,


  \label{e1-baz}


  \end{equation}


где $y_i$, $i=\overline{1,n}$,~--- значения зависимой (объясняемой) 


переменной~$y$; $x_{i1}, x_{i2}, \ldots$\linebreak $\ldots , x_{il}$, $i\hm= \overline{1,n}$,~--- 


значения~$l$~независимых (объясняющих) переменных (регрессоров)~$x_1, 


x_2, \ldots , x_l$; $\varepsilon_i$, $i\hm=\overline{1,n}$,~--- ошибки 


аппроксимации; $\alpha_0, \alpha_1, \ldots , \alpha_l$~--- неизвестные 


параметры; $n$~--- объем выборки.


  


  \textbf{Задача отбора 


<<информативных>> регрессоров.} Необходимо выделить из~$l$~возможных 


регрессоров~$m$~переменных, минимизируя сумму квадратов ошибок для 


регрессии~(1).


  


  В~работах~[7, 8] рассмотрен эффективный алгоритм <<Selection~B>> 


решения поставленной задачи ОИР. В~его основе лежит следующая техника 


оценивания регрессионной модели с~помощью МНК. Предварительно 


проводится нормирование (стандартизация) всех переменных по формулам:


  $$v_i=\fr{y_i-\overline{y}}{\sigma_y}\,;\enskip z_{i1}=\fr{x_{i1}-


\overline{x}_1}{\sigma_{x_1}}\,;\ldots ; z_{il}=\fr{x_{il}-


\overline{x}_l}{\sigma_{x_l}}\,,


  $$


где $\overline{y}, \overline{x}_1, \ldots,  \overline{x}_l$~--- средние значения 


переменных; $\sigma_y, \sigma_{x_1}, \ldots, \sigma_{x_l}$~--- 


сред\-не\-квад\-ра\-тич\-ные отклонения переменных; $v, z_1, \ldots , z_l$~--- 


стандартизованные переменные, для которых среднее значение равно~0, 


а~среднеквадратичное отклонение равно~1.


  


  Тогда регрессии~(1) ставится в~соответствие ее стандартизованная модель:


  \begin{equation}


  v_i=\beta_1 z_{i1}+\beta_2 z_{i2} +\cdots + \beta_l z_{il}+u_i\,,\enskip


  i=\overline{1,n}\,,


  \label{e2-baz}


  \end{equation}


где $\beta_1, \ldots , \beta_l$~--- неизвестные параметры  


(бе\-та-ко\-эф\-фи\-ци\-ен\-ты); $u_i$, $i\hm=\overline{1,n}$,~--- ошиб\-ки 


аппроксимации.


  


  Оценки бета-коэффициентов находятся по формуле:


  \begin{equation*}


  \tilde{\beta}_{\mathrm{МНК}}=r^{-1}_{xx} r_{yx}\,,


  %\label{e3-baz}


  \end{equation*}


где $r_{xx}$~--- матрица коэффициентов парной корреляции между 


объясняющими переменными; $r_{xy}$~--- вектор коэффициентов парной 


корреляции между объясняемой переменной~$y$ и~объясняющими 


переменными $x_1, x_2, \ldots , x_l$.


  


  Коэффициент детерминации для регрессий~(1) и~(2) находится по формуле:


  \begin{equation}


  R^2=r^{\mathrm{T}}_{yx} \tilde{\beta}_{\mathrm{МНК}}\,.


  \label{e4-baz}


  \end{equation}


  


  Алгоритм <<Selection B>> предполагает предварительное задание 


пользователем следующих параметров:


  \begin{enumerate}[1.]


\item Вектор знаков $\mathrm{Sign}_{1\times l}$. Компоненты этого вектора 


назначаются по сле\-ду\-юще\-му правилу:


$$


\hspace*{-4mm}\hspace*{-0.16994pt}\mathrm{Sign}_{1,j}=\begin{cases}


1\,, & \!\!\mbox{если\ $j$-я\ переменная\ оказывает\ положительное\ влияние\ на\ 


$y$};\hspace*{-0.42108pt}\\ 


-1\,, & \!\!\mbox{если\ $j$-я\ переменная\ оказывает\ отрицательное\ влияние\ на\ 


$y$};\\


0\,, & \!\!\mbox{если\ затруднительно\ оценить\ влияние\ $j$-й\ переменной\ на\ 


$y$.}


\end{cases}


$$


  


  Для формирования вектора $\mathrm{Sign}$ желательно привлекать 


экспертов.


\item Число отбираемых регрессоров~$m$.


\item Граница мультиколлинеарности~$d$.


\end{enumerate}





  В алгоритме <<Selection B>> сначала определяется общее число 


альтернативных вариантов моделей $r\hm=C_l^m$, после чего запускается 


основный цикл алгоритма. На каждом повторении формируются необходимые 


матрицы~$r_{xx}$ и~$r_{yx}$, а~затем последовательно проверяются условия:


  \begin{align}


  \left\vert r_{xx}^J\right\vert &<d\,; \label{e5-baz}\\


  \tilde{\beta}_j \mathrm{Sign}_{1,j} &<0\,,\quad j\in J\,,


  \label{e6-baz}


  \end{align}


где $J$~--- множество номеров объясняющих переменных, входящих 


в~регрессию; $r_{xx}^J$~--- матрица парных коэффициентов корреляции для 


объясняющих переменных с~номерами из множества~$J$.


  


  В неравенстве~(\ref{e5-baz}) определитель матрицы парных коэффициентов 


корреляции $\vert r_{xx}^J\vert$ представляет собой критерий обнаружения 


эффекта мультиколлинеарности. Если $\vert r_{xx}^J\vert \hm=0$, то имеет 


место совершенная мультиколлинеарность, а если $\vert r_{xx}^J\vert\hm=1$, то 


мультиколлинеарности нет. Если условие~(\ref{e5-baz}) выполняется, то такая 


регрессия сразу исключается из рассмотрения.


  


  Если выполнено хотя бы одно из условий~(\ref{e6-baz}), то это означает, что 


знак бе\-та-ко\-эф\-фи\-ци\-ен\-та~$\tilde{\beta}_j$ стандартизованной регрессии 


не согласуется с~соответствующим знаком вектора Sign. Следовательно, такая 


модель исключается из дальнейшего <<конкурса>>.


  


  Если очередная альтернативная регрессия прошла <<испытание>> 


условиями~(\ref{e5-baz}) и~(\ref{e6-baz}), то для нее по формуле~(\ref{e4-baz}) 


рассчитывается величина коэффициента детерминации. После окончания 


работы основного цикла определяется наилучшая модель по величине~$R^2$.


  


  Проведенный в~работе~[8] эксперимент показал, что решение задачи ОИР 


с~использование пакета Gretl по алгоритму <<Selection B>> осуществляется 


более чем в~6~раз быстрее, нежели в~ПК АППРМ.


  


  \section{Алгоритм построения вполне интерпретируемых  


и~RTF-адекватных линейных регрессий}


  


  Алгоритм построения вполне интерпретируемых и~RTF-адекват\-ных 


линейных регрессий (см.\ рисунок) был разработан на основе фундаментального 


блока, предложенного в~работе~[7]. Согласно этому блоку, для того чтобы 


очередная альтернативная регрессионная модель была допущена к~участию 


в~<<конкурсе>>, она должна пройти жесткую процедуру отсева, состоящую 


из~5~стадий:


  \begin{enumerate}[(1)]


\item проверка <<информативности>> переменных;


\item проверка присутствия эффекта мультиколлинеарности;


\item проверка соответствия знаков коэффициентов физическому смыслу 


факторов;


\item проверка адекватности по коэффициенту детерминации и~значимости 


в~целом по F-кри\-те\-рию Фишера;


\item проверка значимости коэффициентов по t-кри\-те\-рию Стьюдента.


\end{enumerate}





  Первая стадия связана с~вычислением критерия <<информативности>> 


переменных~\cite{5-baz}:


  \begin{equation*}


  \Gamma\left( x_1, x_2,\ldots , x_l\right)=\sum\limits_{j\in J} \mathrm{Inf}\left( 


x_j\right)\!,


 % \label{e7-baz}


  \end{equation*}


где $\mathrm{Inf}\,(x_j)$~--- некоторый вес (балл) регрессора~$x_j$, например 


в~пятибалльной шкале. Таким образом, этот критерий указывает на 


суммарную <<важность>> вхождения всех независимых переменных 


в~зависимость.


  


  На четвертой стадии для обеспечения значимости регрессии в~целом по 


 \mbox{F-кри}\-те\-рию Фишера значение коэффициента детерминации~$R^2$ не 


должно быть меньше величины


  \begin{equation*}


  R^\prime = \fr{ F_{\mathrm{крит}}(\alpha, m, n-m-1)}{(n-m-


1)/m+F_{\mathrm{крит}}(\alpha, m, n-m-1)}\,,


  %\label{e8-baz}


  \end{equation*}


где $\alpha$~--- заданный уровень значимости; $F_{\mathrm{крит}}(\alpha, 


m, n-m-1)$~--- критическая точка распределения Фишера.





  На пятой стадии вычисление t-кри\-те\-ри\-ев Стюдента для регрессий~(1) 


и~(2) реализуется по формулам:


  \begin{equation*}


  t_j=\fr{\tilde{\beta}_j}{\sqrt{s^2(r_{xx}^J)_{jj}^{-1}}}\,,\quad 


j=\overline{1,m}\,,


 % \label{e9-baz}


  \end{equation*}


 где $s^2\hm= (1\hm-R^2)/ (n\hm-m\hm-1)$~--- величина 


остаточной дисперсии; $(r_{xx}^J)^{-1}_{jj}$~--- $j$-й диагональный элемент матрицы 


$(r_{xx}^J)^{-1}$.


  


  \begin{figure} %fig1


  \vspace*{1pt}


\begin{center}


 \mbox{%


 \epsfxsize=122.7mm 


\epsfbox{baz-1.eps}


 }





\vspace*{4pt}





  {\small Алгоритм построения линейных регрессий}


     \end{center}


     \vspace*{-6.72816pt}


  \end{figure}


  


\pagebreak


  


  Пользователю предварительно необходимо задать следующие параметры:


  \begin{enumerate}[1.]


\item Вектор знаков $\mathrm{Sign}_{1\times l}$.


\item Вектор <<информативности>> переменных $\mathrm{Inf}_{1\times l}$. 


Компоненты этого вектора представляют собой целые числа из 


интервала~[1,\,5]. Если $\mathrm{Inf}_{1\times j}\hm=5$, то <<информативность>> 


\mbox{$j$-й} переменной максимальна и~составляет 5~баллов, а~если 


$\mathrm{Inf}_{1\times j}\hm=1$, то <<информативность>> минимальна 


и~составляет 1~балл. Для формирования этого вектора также желательно 


привлекать экспертов.


\item Число отбираемых регрессоров~$m$.


\item Границу мультиколлинеарности~$d$.


\item Границу значимости для F-кри\-те\-рия Фишера~$R^\prime$.


\item Границу значимости для t-кри\-те\-ри\-ев Стьюдента~$t^\prime$.


\end{enumerate}


  


  На рисунке $\mathrm{Glavn}_{r\times m}$~--- матрица альтернатив, в~которой 


объясняемая переменная имеет номер~1, а объясняющие переменные~--- $2, 3, 


\ldots , l\hm+1$; adeq~--- счетчик вполне интерпретируемых  


и~RTF-адекват\-ных моделей; flag~--- бинарная переменная, которая отвечает 


за выполнение всех пяти требований, предъявляемым к~моделям; Krit~--- 


критериальная матрица, в~первом столбце которой содержатся значения 


коэффициента детерминации~$R^2$, во втором~--- определителя $\vert 


r_{xx}\vert$, в~третьем~--- критерия <<информативности>> переменных, 


в~остальных~--- номера объясняющих переменных.


  


  Предложенный алгоритм был реализован в~виде программы 


с~использованием эконометрического пакета Gretl.


  


  \section{Заключение}


  


  Разработанная программа построения вполне интерпретируемых  


и~RTF-аде\-кват\-ных линейных регрессионных моделей позволяет находить 


среди многочисленных альтернатив по-настоящему качественные регрессии, 


значимые в~целом по F-кри\-те\-рию Фишера, в~которых все знаки 


коэффициентов соответствуют физическому смыслу факторов, незначителен 


эффект мультиколлинеарности и~все коэффициенты значимы по t-кри\-те\-рию 


Стьюдента. Кроме того, выбор состава переменных приобретает определенную 


сознательность, поскольку в~процесс построения модели привлекаются 


эксперты в~исследуемой предметной об\-ласти.


  


{\small\frenchspacing


{%\baselineskip=10.8pt


\addcontentsline{toc}{section}{Литература}


\begin{thebibliography}{9}


\bibitem{1-baz}


\Au{Pardoe I.} Applied regression modeling.~--- Hoboken, NJ, USA: Wiley, 2020. 336~p.


\bibitem{2-baz}


\Au{Miller A.\,J.} Subset selection in regression.~--- London, U.K.: Chapman \& Hall/CRC, 2002. 


256~p.


\bibitem{3-baz}


\Au{Venkatesh B., Anuradha~J.} A~review of feature selection and its methods~/\!\!/ Cybernetics 


Information Technologies, 2019. Vol.~19. P.~3--26.


\bibitem{4-baz}


\Au{Стрижов В.\,В., Крымова Е.\,А.} Методы выбора регрессионных моделей.~--- М.:  


ВЦ РАН, 2010. 60~с.





\bibitem{6-baz}


\Au{Носков С.\,И., Потороченко Н.\,А.} Диалоговая система реализации <<конкурса>> 


регрессионных зависимостей~/\!\!/ Управ\-ля\-ющие сис\-те\-мы и~машины, 1992. №\,3-4. 


С.~111--116.


\bibitem{5-baz}


\Au{Носков С.\,И.} Технология моделирования объектов с~нестабильным 


функционированием и~неопределенностью в~данных.~--- Иркутск: Облинформпечать, 1996. 


321~c.





\bibitem{7-baz}


\Au{Базилевский М.\,П.} Фундаментальный блок алгоритмов построения хорошо 


интерпретируемых качественных регрессионных моделей~/\!\!/ Информационные 


технологии и~математическое моделирование в~управлении сложными системами, 2020. 


 №\,3(8). С.~1--10.


\bibitem{8-baz}


\Au{Базилевский М.\,П.} Повышение эффективности алгоритма отбора по критерию 


детерминации информативных регрессоров в~регрессионных моделях~/\!\!/ Прикладная 


математика и~информатика: современные исследования в~области естественных 


и~технических наук.~--- Тольятти, 2018. С.~196--202.





    \end{thebibliography}


} }











\vspace*{-3pt}





\hfill{\small\textit{Поступила в~редакцию 13.01.21}}





\vspace*{8pt}





\hrule





\vspace*{2pt}





%\newpage





%\vspace*{-28pt}





\hrule





%\vspace*{8pt}





\def\tit{A PROGRAM FOR~CONSTRUCTING OF~QUITE INTERPRETABLE AND~RTF-ADEQUATE 


LINEAR REGRESSION MODELS}








\def\titkol{A program for~constructing of~quite interpretable and~rtf-adequate 


linear regression models}





\def\aut{M.\,P.~Bazilevskiy}





\def\autkol{M.\,P.~Bazilevskiy}





\titel{\tit}{\aut}{\autkol}{\titkol}





\vspace*{-15pt}





\noindent


Department of Mathematics, Irkutsk State Transport University, 15~Chernyshevskogo Str., Irkutsk 664074, 


Russian Federation








\def\leftfootline{\small{\textbf{\thepage}


\hfill Sistemy i Sredstva Informatiki~---


Systems and Means of Informatics 2021 vol~31 no\,4}


}%


 \def\rightfootline{\small{Sistemy i Sredstva Informatiki~---


 Systems and Means of Informatics 2021 vol~31 no\,4


\hfill \textbf{\thepage}}}





\vspace*{3pt}





\Abste{The article is devoted to the problem of feature selection in regression models estimated using the 


ordinary least squares method. Models constructed as a~result of such selection are often inadequate and 


poorly interpreted. For the first time, the definitions of ``quite interpretable'' and ``RTF-adequate''


 regression 


models are formulated. The previously proposed effective algorithm for solving the problem of feature 


selection is considered. On its basis, an algorithm has been developed for constructing quite interpretable and 


RTF-adequate linear regression models. In it, for each regression, the following tests are sequentially carried 


out: ``informativeness'' of variables, multicollinearity, correspondence of coefficients signs to the physical 


meaning of factors, adequacy of model in terms of coefficient of determination and significance in general 


according to Fisher's F-test, and significance of the coefficients according to the Student's t-test. The proposed 


algorithm is implemented as a~program for the Gretl econometric package. The developed program is 


universal and can be used to solve a~wide range of data analysis tasks.}





\KWE{feature selection; ordinary least squares; quite interpretable and RTF-adequate regression; 


variable ``informativeness'' criterion; multicollinearity; Fisher's F-test; Student's t-test}





\DOI{10.14357\!/\!08696527210402} 





%\vspace*{-16pt}





%\Ack


%\noindent








%\vspace*{-6pt}





\renewcommand{\bibname}{\protect\rmfamily References}


%\renewcommand{\bibname}{\large\protect\rm References}





{\small\frenchspacing


{%\baselineskip=10.8pt


\addcontentsline{toc}{section}{References}


\begin{thebibliography}{9}


\bibitem{1-baz-1}


\Aue{Pardoe, I.} 2020. \textit{Applied regression modeling}. Hoboken, NJ, USA: Wiley. 336 p.


\bibitem{2-baz-1}


\Aue{Miller, A.\,J.} 2002. \textit{Subset selection in regression}. London, U.K.: Chapman \& Hall/CRC. 256 


p.


\bibitem{3-baz-1}


\Aue{Venkatesh, B., and J.~Anuradha.} 2019. A~review of feature selection and its methods. 


\textit{Cybernetics Information Technologies} 19:3--26.


\bibitem{4-baz-1}


\Aue{Strizhov, V.\,V., and E.\,A.~Krymova.} 2010. \textit{Metody vybora regressionnykh modeley} 


[Regression model selection methods]. Moscow: CC RAS. 60~p.





\bibitem{6-baz-1}


\Aue{Noskov, S.\,I., and N.\,A.~Potorochenko.} 1992. Dialogovaya sistema realizatsii konkursa 


regressionnykh zavisimostey [Dialogue system for the implementation of the competition of regression 


dependencies]. \textit{Upravlyayushchie sistemy i~mashiny} [Control Systems and Computers]  


3-4:111--116.


\bibitem{5-baz-1}


\Aue{Noskov, S.\,I.} 1996. \textit{Tekhnologiya modelirovaniya ob''ektov s~nestabil'nym 


funk\-tsi\-o\-ni\-ro\-va\-ni\-em i~neopredelennost'yu v~dannykh} [Technology for modeling objects with unstable 


functioning and uncertainty in data]. Irkutsk: Oblinformpechat'. 321~p.





\bibitem{7-baz-1}


\Aue{Bazilevskiy, M.\,P.} 2020. Fundamental'nyy blok algoritmov postroeniya khorosho 


interpretiruemykh kachestvennykh regressionnykh modeley [The fundamental block of algorithms for 


constructing well-interpreted qualitative regression models]. \textit{Informatsionnye tekhnologii 


i~matematicheskoe modelirovanie v~upravlenii slozhnymi sistemami} [Information Technology and 


Mathematical Modeling in the Management of Complex Systems] 3(8):1--10.


\bibitem{8-baz-1}


\Aue{Bazilevskiy, M.\,P.} 2018. Povyshenie effektivnosti algoritma otbora po kriteriyu determinatsii 


informativnykh regressorov v~regressionnykh modelyakh [Improving the efficiency of the algorithm for 


selecting informative regressors by the criterion of determination in regression models]. 


\textit{Prikladnaya matematika i~informatika: sovremennye issledovaniya v~oblasti estestvennykh


i~tekhnicheskikh nauk}


[Applied Mathematics and Informatics: Contemporary Research in Natural 


and Technical Sciences]. Tol'yatti. 196--202.





\end{thebibliography}


} }





\vspace*{-3pt}





\hfill{\small\textit{Received January 13, 2021}} 





\vspace*{-14pt}  





\Contrl








\noindent


\textbf{Bazilevskiy Mikhail P.} (b.\ 1987)~--- Candidate of Science (PhD) in technology, associate 


professor, Irkutsk State Transport University, 15~Chernyshevkogo Str., Irkutsk 664074, Russian Federation; 


\mbox{mik2178@yandex.ru}





\label{end\stat}





\renewcommand{\bibname}{\protect\rm Литература} 


